\chapterbackmatter{Supplementary Exercises}

\begin{enumerate}
  \SupplementaryExercise{1}{Grassmann Delta Saturation}
    {(Adapted from Gates, Grisaru, Rocek, and Siegel, \emph{Superspace},
    Chapter 6, Section 6.3)}
    {gates-grisaru-rocek-siegel-superspace-ch30}
  Starting from Eq.~\eqref{eq:26.6.8}, prove that a chain of \(r\)
  superspace propagators on a tree identifies all Grassmann coordinates
  and leaves one \(d^4\theta\) integral.  Explain why an additional
  undifferentiated delta function closing a loop makes the result vanish.

  \SupplementaryExercise{2}{Chiral Projection Test}
    {(Adapted from Gates, Grisaru, Rocek, and Siegel, \emph{Superspace},
    Chapter 6, Sections 6.1 and 6.3)}
    {gates-grisaru-rocek-siegel-superspace-ch30}
  For
  \(\mathcal P_L=-\mathcal D_L^2\mathcal D_R^2/(16\Box)\), verify
  \(\mathcal P_L^2=\mathcal P_L\) and
  \(\mathcal D_R\mathcal P_L=0\).  State precisely which zero-momentum
  modes must be excluded when writing \(\Box^{-1}\).

  \SupplementaryExercise{3}{Potential-Field Redundancy}
    {(Adapted from Gates, Grisaru, Rocek, and Siegel, \emph{Superspace},
    Chapter 6, Section 6.1)}
    {gates-grisaru-rocek-siegel-superspace-ch30}
  Let \(\Phi=\mathcal D_R^2S\).  Show directly that
  \(S\to S+\mathcal D_RX\) leaves \(\Phi\) unchanged, and prove that
  coupling \(S\) to a source is invariant only when the source obeys the
  corresponding superspace transversality condition.

  \SupplementaryExercise{4}{Projector Completeness}
    {(Adapted from Gates, Grisaru, Rocek, and Siegel, \emph{Superspace},
    Chapter 6, Sections 6.2--6.3)}
    {gates-grisaru-rocek-siegel-superspace-ch30}
  Define \(\mathcal P_R=-\mathcal D_R^2\mathcal D_L^2/(16\Box)\) and
  \(\mathcal P_T=\mathbf 1-\mathcal P_L-\mathcal P_R\).  Establish the
  orthogonality and completeness relations for these three projectors
  and identify the gauge-invariant sector of a real vector superfield.

  \SupplementaryExercise{5}{Weighted Vector Gauge}
    {(Adapted from Gates, Grisaru, Rocek, and Siegel, \emph{Superspace},
    Chapter 6, Section 6.2)}
    {gates-grisaru-rocek-siegel-superspace-ch30}
  Suppose gauge fixing changes the Abelian vector quadratic operator to
  \(-\Box[\mathcal P_T+\alpha^{-1}(\mathcal P_L+\mathcal P_R)]\).
  Invert it and show that every dependence on \(\alpha\) disappears
  between transverse external sources.

  \SupplementaryExercise{6}{Massive Chiral Kernel}
    {(Adapted from Gates, Grisaru, Rocek, and Siegel, \emph{Superspace},
    Chapter 6, Section 6.3)}
    {gates-grisaru-rocek-siegel-superspace-ch30}
  Add \(m[\Phi^2]_{\mathcal F}/2+\mathrm{H.c.}\) to a free chiral
  theory.  Invert the quadratic \(2\times2\) chiral--antichiral kernel
  and determine the pole shared by all four component propagators.

  \SupplementaryExercise{7}{Derivative Transfer Rule}
    {(Adapted from Gates, Grisaru, Rocek, and Siegel, \emph{Superspace},
    Chapter 6, Section 6.3.e)}
    {gates-grisaru-rocek-siegel-superspace-ch30}
  Prove the superspace integration-by-parts rule that transfers
  \(\mathcal D_{L,R}\) from one end of a propagator to the other.
  Include the Grassmann parity sign and apply the result to a product
  of two fermionic superfields.

  \SupplementaryExercise{8}{One-Loop Self-Energy Algebra}
    {(Adapted from Gates, Grisaru, Rocek, and Siegel, \emph{Superspace},
    Chapter 6, Sections 6.3.e and 6.4)}
    {gates-grisaru-rocek-siegel-superspace-ch30}
  In the massless cubic chiral theory, perform only the \(D\)-algebra
  of the one-loop two-point supergraph.  Show that its local
  divergence is proportional to a full-superspace integral
  \([\Phi^*\Phi]_D\), not an \(\mathcal F\)-term.

  \SupplementaryExercise{9}{Vanishing Vacuum Loop}
    {(Adapted from Gates, Grisaru, Rocek, and Siegel, \emph{Superspace},
    Chapter 6, Sections 6.3--6.4)}
    {gates-grisaru-rocek-siegel-superspace-ch30}
  Consider a connected vacuum supergraph made only of massless chiral
  propagators and superpotential vertices.  Use unsaturated Grassmann
  delta functions, rather than component cancellations, to show that
  its contribution to the vacuum energy vanishes.

  \SupplementaryExercise{10}{Superficial Divergence Bound}
    {(Adapted from Gates, Grisaru, Rocek, and Siegel, \emph{Superspace},
    Chapter 6, Section 6.3.d)}
    {gates-grisaru-rocek-siegel-superspace-ch30}
  For a renormalizable chiral supergraph with \(L\) loops, \(I\)
  internal lines, and \(V\) vertices, combine momentum counting with
  the \(D\)-algebra to show why divergent terms require at most two
  external chiral superfields.

  \SupplementaryExercise{11}{One Overall Superspace Point}
    {(Adapted from Gates, Grisaru, Rocek, and Siegel, \emph{Superspace},
    Chapter 6, Section 6.3.c)}
    {gates-grisaru-rocek-siegel-superspace-ch30}
  Give an inductive proof that every connected supergraph reduces to
  one overall Grassmann coordinate after its \(D\)-algebra is
  completed.  Explain why this establishes superspace locality without
  implying locality in the spacetime coordinates.

  \SupplementaryExercise{12}{Cubic Vertex Counting}
    {(Adapted from Gates, Grisaru, Rocek, and Siegel, \emph{Superspace},
    Chapter 6, Sections 6.3.a--6.3.b)}
    {gates-grisaru-rocek-siegel-superspace-ch30}
  A connected graph in a cubic chiral theory has \(V\) chiral vertices,
  \(V^*\) antichiral vertices, \(I\) internal lines, and \(E,E^*\)
  external lines.  Derive the incidence relations and use them to
  reproduce the loop formula in Section~30.3.

  \SupplementaryExercise{13}{Yukawa Beta Tensor}
    {(Adapted from Antusch and Ratz, \emph{Supergraph Techniques and
    Two-Loop Beta-Functions}, Sections 2.1 and 3)}
    {antusch-ratz-supergraphs-2002-ch30}
  Let \(W=Y_{ijk}\Phi_i\Phi_j\Phi_k/6\), with \(Y_{ijk}\) symmetric.
  Assuming the superpotential has no independent counterterm, derive
  its beta function in terms of the anomalous-dimension matrix
  \(\gamma_i{}^j\).

  \SupplementaryExercise{14}{Mass-Matrix Running}
    {(Adapted from Antusch and Ratz, \emph{Supergraph Techniques and
    Two-Loop Beta-Functions}, Sections 2.1 and 3.3)}
    {antusch-ratz-supergraphs-2002-ch30}
  For \(W=M_{ij}\Phi_i\Phi_j/2\), derive the renormalization-group
  equation for the symmetric mass matrix \(M\).  Check that the answer
  transforms covariantly under a constant unitary change of chiral
  field basis.

  \SupplementaryExercise{15}{Higher-Operator Running}
    {(Adapted from Antusch and Ratz, \emph{Supergraph Techniques and
    Two-Loop Beta-Functions}, Sections 2.1 and 3.2)}
    {antusch-ratz-supergraphs-2002-ch30}
  A superpotential contains
  \(C_{i_1\cdots i_n}\Phi_{i_1}\cdots\Phi_{i_n}/(n!\,M^{n-3})\).
  Derive the beta function of \(C\) from wavefunction renormalization
  and explain why the result remains valid for \(n>3\).

  \SupplementaryExercise{16}{One-Loop Anomalous Dimension}
    {(Adapted from Antusch and Ratz, \emph{Supergraph Techniques and
    Two-Loop Beta-Functions}, Section 2.1)}
    {antusch-ratz-supergraphs-2002-ch30}
  For chiral fields with cubic couplings \(Y_{ijk}\) and gauge
  coupling \(g\), show that the one-loop anomalous dimension must have
  the tensor form
  \(aY_{imn}Y^{*jmn}-bg^2C_2(R_i)\delta_i{}^j\).
  Determine the signs from physical wavefunction screening.

  \SupplementaryExercise{17}{Casimir Conversion}
    {(Adapted from Antusch and Ratz, \emph{Supergraph Techniques and
    Two-Loop Beta-Functions}, Eqs. (2.6)--(2.7))}
    {antusch-ratz-supergraphs-2002-ch30}
  Prove
  \(C_2(R)\dim R=T(R)\dim G\) by tracing the defining Casimir equation.
  Evaluate \(C_2\) and \(T\) for the fundamental representation of
  \(SU(N)\) with \(T(\mathbf N)=1/2\).

  \SupplementaryExercise{18}{Abelian Charge Factor}
    {(Adapted from Antusch and Ratz, \emph{Supergraph Techniques and
    Two-Loop Beta-Functions}, Section 2.1)}
    {antusch-ratz-supergraphs-2002-ch30}
  Replace a non-Abelian gauge factor by \(U(1)\).  Demonstrate from the
  two gauge vertices on a chiral self-energy line that
  \(C_2(R)\) is replaced by \(q^2\), and explain why no term linear in
  \(q\) can survive.

  \SupplementaryExercise{19}{Field-Redefinition Contribution}
    {(Adapted from Antusch and Ratz, \emph{Supergraph Techniques and
    Two-Loop Beta-Functions}, Sections 2.1 and 3)}
    {antusch-ratz-supergraphs-2002-ch30}
  Let \(\Phi_B=Z^{1/2}\Phi\) with matrix-valued \(Z\).  Differentiate
  this relation at fixed bare field and derive the contribution of
  each external chiral leg to the running of an arbitrary
  superpotential tensor.

  \SupplementaryExercise{20}{EFT Triangularity}
    {(Adapted from Antusch and Ratz, \emph{Supergraph Techniques and
    Two-Loop Beta-Functions}, Sections 1--2.1)}
    {antusch-ratz-supergraphs-2002-ch30}
  Explain by dimensional analysis why an operator suppressed by
  \(M^{-k}\) does not affect renormalizable wavefunction counterterms
  at leading order in \(1/M\), while renormalizable interactions can
  renormalize that higher-dimensional operator.

  \SupplementaryExercise{21}{Holomorphic versus Canonical Coupling}
    {(Adapted from Antusch and Ratz, \emph{Supergraph Techniques and
    Two-Loop Beta-Functions}, Sections 1--2.1)}
    {antusch-ratz-supergraphs-2002-ch30}
  A holomorphic cubic coefficient is scale independent before fields
  are canonically normalized.  Show explicitly how a nonzero physical
  beta function nevertheless arises when
  \(\Phi_{\mathrm c}=Z^{1/2}\Phi\).

  \SupplementaryExercise{22}{Background Hessian}
    {(Adapted from Kuzenko and Tyler, \emph{The One-Loop Effective
    Potential of the Wess-Zumino Model Revisited}, Section 2)}
    {kuzenko-tyler-wz-effective-action-2014-ch30}
  Expand a general Wess--Zumino action about a chiral background
  \(\Phi\).  Construct the quadratic chiral--antichiral Hessian in
  terms of \(\Psi=W''(\Phi)\), and show why cubic quantum vertices
  cannot contribute at one loop.

  \SupplementaryExercise{23}{Logarithm of the Hessian}
    {(Adapted from Kuzenko and Tyler, \emph{The One-Loop Effective
    Potential of the Wess-Zumino Model Revisited}, Section 2.2)}
    {kuzenko-tyler-wz-effective-action-2014-ch30}
  Evaluate the Gaussian functional integral for the background
  Hessian and obtain
  \(\Gamma^{(1)}=(i/2)\operatorname{Tr}\ln(H/H_0)\).
  Expand the logarithm and identify the symmetry factor of an
  \(n\)-vertex ring supergraph.

  \SupplementaryExercise{24}{Kahler Dimensional Check}
    {(Adapted from Kuzenko and Tyler, \emph{The One-Loop Effective
    Potential of the Wess-Zumino Model Revisited}, Sections 1 and 3.1)}
    {kuzenko-tyler-wz-effective-action-2014-ch30}
  In four dimensions \([\Phi]=1\) and \([\Psi=W''(\Phi)]=1\).
  Use dimensions and phase invariance to determine the possible
  one-loop local Kahler correction for constant \(\Psi\), up to its
  overall coefficient and a local subtraction.

  \SupplementaryExercise{25}{Logarithmic Momentum Integral}
    {(Adapted from Kuzenko and Tyler, \emph{The One-Loop Effective
    Potential of the Wess-Zumino Model Revisited}, Sections 3.1 and 4.1)}
    {kuzenko-tyler-wz-effective-action-2014-ch30}
  Differentiate
  \(J(M^2)=\int d^dp\,(2\pi)^{-d}\ln(p^2+M^2)\) with respect to
  \(M^2\), evaluate the pole in dimensional regularization, and
  reconstruct the \(M^2\ln(M^2/\mu^2)\) dependence.

  \SupplementaryExercise{26}{Four-Derivative Invariant}
    {(Adapted from Kuzenko and Tyler, \emph{The One-Loop Effective
    Potential of the Wess-Zumino Model Revisited}, Sections 1 and 3.2)}
    {kuzenko-tyler-wz-effective-action-2014-ch30}
  For a constant chiral background, build the lowest real scalar
  full-superspace integrand that vanishes when all spinor derivatives
  of \(\Phi\) vanish.  Show that it begins with two left and two right
  superderivatives.

  \SupplementaryExercise{27}{Auxiliary-Field Order}
    {(Adapted from Kuzenko and Tyler, \emph{The One-Loop Effective
    Potential of the Wess-Zumino Model Revisited}, Sections 1 and 5)}
    {kuzenko-tyler-wz-effective-action-2014-ch30}
  Set spacetime derivatives and fermions to zero in the invariant of
  Exercise S.30.26.  Determine its dependence on the scalar and
  auxiliary components and explain why it belongs to an auxiliary
  field potential rather than the ordinary Kahler potential.

  \SupplementaryExercise{28}{Massless Scaling Form}
    {(Adapted from Kuzenko and Tyler, \emph{The One-Loop Effective
    Potential of the Wess-Zumino Model Revisited}, Sections 1 and 6)}
    {kuzenko-tyler-wz-effective-action-2014-ch30}
  In the massless cubic model, use dimensions and constant phase
  transformations of \(\Phi\) to find the denominator required beneath
  the four-superderivative numerator of Exercise S.30.26.  Verify that
  the resulting integrand is dimension two.

  \SupplementaryExercise{29}{Heavy-Field Decoupling}
    {(Adapted from Kuzenko and Tyler, \emph{The One-Loop Effective
    Potential of the Wess-Zumino Model Revisited}, Sections 1 and 3)}
    {kuzenko-tyler-wz-effective-action-2014-ch30}
  For \(W''(\Phi)=m+\lambda\Phi\) with
  \(|\lambda\Phi|\ll |m|\), expand the one-loop Kahler correction
  through quadratic order in \(\Phi\).  Identify the removable linear
  term and the threshold correction to the kinetic term.

  \SupplementaryExercise{30}{Regulator Ordering Warning}
    {(Adapted from Kuzenko and Tyler, \emph{The One-Loop Effective
    Potential of the Wess-Zumino Model Revisited}, Section 4.2)}
    {kuzenko-tyler-wz-effective-action-2014-ch30}
  Consider a conditionally convergent proper-time and momentum
  integral arising after \(D\)-algebra.  Explain why exchanging the
  order of integration can change a finite coefficient, and formulate
  a regulator-first prescription that removes the ambiguity.
\end{enumerate}
